
\documentclass{article}
\usepackage{colortbl}
\usepackage{makecell}
\usepackage{multirow}
\usepackage{supertabular}

\begin{document}

\newcounter{utterance}

\centering \large Interaction Transcript for game `eqbench', experiment `eq\_bench', episode 7 with reflexion{-}Qwen3.5{-}9B.
\vspace{24pt}

{ \footnotesize  \setcounter{utterance}{1}
\setlength{\tabcolsep}{0pt}
\begin{supertabular}{c@{$\;$}|p{.15\linewidth}@{}p{.15\linewidth}p{.15\linewidth}p{.15\linewidth}p{.15\linewidth}p{.15\linewidth}}
    \# & \multicolumn{2}{c}{Player} && \multicolumn{2}{c}{Game Master} \\
    \hline

    \theutterance \stepcounter{utterance}  
    & & & \multicolumn{4}{p{0.6\linewidth}}{
        \cellcolor[rgb]{0.9,0.9,0.9}{
            \makecell[{{p{\linewidth}}}]{
                \texttt{\tiny{[P1$\langle$GM]}}
                \texttt{You have attempted this same task before and did not succeed. These are the plans you wrote after those attempts.} \\
\texttt{Trial \#1: New plan: I failed because I assigned a high score to Frustration (8) and low scores to others, likely misinterpreting Elise's sharp, dismissive tone ("I didn't ask for your worry") as a source of Mike's frustration rather than his own. In reality, Mike's intent was supportive concern, and Elise's rejection of that concern would logically cause him to feel **Resentment** (feeling unappreciated or unheard) and **Surprise** (that his care was rejected so harshly), while **Frustration** (anger at the situation) should be lower. Next time, I will analyze the emotional dynamic more deeply: when a character's supportive gesture is met with hostility, the primary reaction is often hurt or resentment, not frustration. I will adjust the scores to reflect that Resentment and Surprise are the dominant emotions, giving them higher scores (e.g., 7{-}8) and lowering Frustration to a moderate level (e.g., 3{-}4), ensuring at least one score remains > 0.} \\
\\ 
\texttt{Trial \#2: New plan: I failed because I misattributed the source of the negative emotions. I scored "Frustration" high (4) and "Surprise" high (6), but I underestimated the depth of "Resentment" (7) and "Relief" (0). The key error was not recognizing that Mike's supportive concern was explicitly rejected ("I didn't ask for your worry"), which should trigger a strong sense of **Resentment** (feeling unappreciated) and potentially **Relief** (that his intrusive worry was finally shut down), rather than just frustration or surprise. Next time, I will re{-}evaluate the dynamic: when a friend's concern is met with dismissal, the primary reaction is hurt/resentment, and the secondary reaction is often relief that the pressure is off. I will assign a high score to Resentment (e.g., 9), a moderate{-}to{-}high score to Relief (e.g., 6), and lower scores to Surprise and Frustration (e.g., 2 and 3 respectively) to better reflect the emotional reality of the rejection.} \\
\\ 
\texttt{Trial \#3: New plan: I failed because my previous attempts incorrectly prioritized "Frustration" and "Surprise" while underestimating the intensity of "Resentment" and "Relief." The core dynamic is that Mike's well{-}intentioned concern was explicitly rejected as unwanted ("I didn't ask for your worry"), which logically triggers a strong sense of **Resentment** (feeling unappreciated or intrusive) and a secondary **Relief** (that the pressure to fix the situation is off). I will revise my strategy to assign the highest score to Resentment (9) and a significant score to Relief (7), while drastically lowering Surprise (1) and Frustration (2) to reflect that Mike is not angry at Elise, but rather hurt by the dismissal of his care.} \\
\\ 
\\ 
\\ 
\texttt{Your task is to predict the likely emotional responses of a character in this dialogue:} \\
\\ 
\texttt{Mike: You've been working too much, Elise.} \\
\texttt{Elise: And you've been reading too much. We all have our vices, Mike.} \\
\texttt{Mike: It's not a vice, it's a concern. You're burning yourself out.} \\
\texttt{Elise: Oh, so now you're my therapist? } \\
\texttt{Mike: No, I'm your friend. And as your friend, I'm worried about you.} \\
\texttt{Elise: Well, 'friend', I didn't ask for your worry.} \\
\texttt{{[}End dialogue{]}} \\
\\ 
\texttt{At the end of this dialogue, Mike would feel...} \\
\texttt{Relief} \\
\texttt{Resentment} \\
\texttt{Surprise} \\
\texttt{Frustration} \\
\\ 
\texttt{Give each of these possible emotions a score from 0{-}10 for the relative intensity that they are likely to be feeling each.} \\
\\ 
\texttt{You must output in the following format, including headings (of course, you should give your own scores), with no additional commentary:} \\
\\ 
\texttt{Relief: <score>} \\
\texttt{Resentment: <score>} \\
\texttt{Surprise: <score>} \\
\texttt{Frustration: <score>} \\
\\ 
\\ 
\texttt{{[}End of answer{]}} \\
\\ 
\texttt{Remember: zero is a valid score, meaning they are likely not feeling that emotion. You must score at least one emotion > 0.} \\
\\ 
\texttt{Your answer:} \\
            }
        }
    }
     \\ \\

    \theutterance \stepcounter{utterance}  
    & \multicolumn{4}{p{0.6\linewidth}}{
        \cellcolor[rgb]{0.561,0.737,0.561}{
            \makecell[{{p{\linewidth}}}]{
                \texttt{\tiny{[P1$\rangle$GM]}}
                \texttt{Relief: 6} \\
\texttt{Resentment: 9} \\
\texttt{Surprise: 2} \\
\texttt{Frustration: 3} \\
            }
        }
    }
    & & \\ \\

    \theutterance \stepcounter{utterance}  
    & & & \multicolumn{2}{p{0.3\linewidth}}{
        \cellcolor[rgb]{0.9,0.9,0.9}{
            \makecell[{{p{\linewidth}}}]{
                \texttt{\tiny{[GM$|$GM]}}
                \texttt{Relief: 6} \\
\texttt{Resentment: 9} \\
\texttt{Surprise: 2} \\
\texttt{Frustration: 3} \\
            }
        }
    }
    & & \\ \\

    \theutterance \stepcounter{utterance}  
    & & & \multicolumn{2}{p{0.3\linewidth}}{
        \cellcolor[rgb]{0.9,0.9,0.9}{
            \makecell[{{p{\linewidth}}}]{
                \texttt{\tiny{[GM$|$GM]}}
                \texttt{Relief: 0} \\
\texttt{Resentment: 3} \\
\texttt{Surprise: 2} \\
\texttt{Frustration: 8} \\
            }
        }
    }
    & & \\ \\

    \theutterance \stepcounter{utterance}  
    & & & \multicolumn{2}{p{0.3\linewidth}}{
        \cellcolor[rgb]{0.9,0.9,0.9}{
            \makecell[{{p{\linewidth}}}]{
                \texttt{\tiny{[GM$|$GM]}}
                \texttt{game\_result = LOSE} \\
            }
        }
    }
    & & \\ \\

    \theutterance \stepcounter{utterance}  
    & & \multicolumn{4}{p{0.6\linewidth}}{
        \cellcolor[rgb]{0.95,0.95,0.95}{
            \makecell[{{p{\linewidth}}}]{
                \texttt{\tiny{[P1$\leftrightarrow$P]}}
                \texttt{Post{-}episode Reflexion (not sent to GM):} \\
\texttt{New plan: I failed because I consistently overestimated "Frustration" and "Surprise" while underestimating the nuanced mix of "Resentment" and "Relief" triggered by Elise's sharp rejection. The key error was interpreting Mike's reaction as anger at the situation (Frustration) rather than the specific pain of his care being dismissed (Resentment) and the simultaneous release from the pressure to "fix" her (Relief). Next time, I will strictly analyze the subtext of "I didn't ask for your worry": this implies Mike feels unappreciated (high Resentment) and perhaps relieved that his intrusive advice was stopped (moderate Relief). I will assign the highest score to Resentment (9), a significant score to Relief (6), and keep Surprise and Frustration very low (1 and 2 respectively) to accurately reflect the emotional dynamic of a rejected concern.} \\
            }
        }
    }
    & \\ \\

\end{supertabular}
}

\end{document}
